\documentclass[11pt]{article}
\usepackage{amsmath,amssymb}
\title{Room 14 (seat: direct attempt): route (b) done properly ---
a modulus-only bound on the resonance-pinned $N=6$ legs}
\date{2026-09-27}
\begin{document}
\maketitle

\section*{0. What was wrong with the earlier probe}

\texttt{room13\_status.tex} flagged its own five-minute check as unreliable because it tested
$w=ce^{-x}$ (some unspecified $c$, no factor of $2$ on $x$) instead of the lemma's actual
resonance-relative quantity. This room re-derives that quantity from \texttt{lem:qi-bounds}'s proof
text itself and redoes the computation with it.

\section*{1. The exact definition of $w$, quoted from the source}

From the proof of \texttt{lem:qi-bounds} (\emph{A ray lemma} paragraph), \texttt{paper2a.tex}:
\begin{quote}
``Let $|c|=1$, $G(v)=\log(1-ce^{-v})$, $|ce^{-v_0}|\le1$, and
$\delta=\inf_{\rho\ge0}|1-ce^{-v_0-\rho e^{i\theta}}|>0$. \ldots Along the ray $w=ce^{-v}$ has
$|w|\le e^{-a'-\rho\cos\theta}$ and $|1-w|\ge\max(1-|w|,\sigma(\arg w))$\ldots''
\end{quote}
So $w(\rho)=c\,e^{-v_0-\rho e^{i\theta}}$, and $|c|=1$ is a specific root of unity fixed by which
factor of the theta-function product is being bounded. From Part (i)'s proof: ``For $f_e$,
$\log(i^je^{-2x-jt};P)_\infty=\sum_mG_j(2x+jt+4tm)$ with $c=i^j$''\,---\,i.e.\ $v=2x+jt+4tm$, so the
$x$-dependence of $v$ enters as $2x$, not $x$ (this is exactly the factor Room 13's probe dropped).
From Part (ii)'s proof: ``With $v=2x+jt-\frac{i\pi\gamma_j}2$ ($c_j=e^{i\pi\gamma_j/2}$)\ldots The
factors $(P/z_j;P)_\infty$ have argument $c_j^{-1}e^{2x}e^{-(4-j)t}$ of modulus $\le1$ and are
treated by the ray lemma with $\delta\ge d'(x)$.'' So the resonance constant is one of the roots of
unity $c_j=e^{i\pi\gamma_j/2}$ (for $N=4$: 8th roots, indexed by residue class $j$), and, dropping
the $O(t)$ shift (which is what vanishes in the $t\to0$ asymptotic regime the whole construction
lives in), the leading-order resonance-relative quantity is
\[
   w(x)=c_j\,e^{-2x},\qquad |c_j|=1.
\]
This is the quantity used throughout; Room 13's probe used $w=ce^{-x}$ (wrong power of $x$, and an
unspecified $c$), which is exactly why it was flagged unreliable.

\section*{2. The $N=6$ analogue: which resonance constant}

\texttt{room8\_seat\_hilbert.tex} identifies the five $N=6$ resonance loci as
$\operatorname{Im}(x)\equiv\arg(\zeta_6^i)/2\pmod\pi$, $i=1,\ldots,5$. The $N$-generic Pochhammer
decomposition $(q;q)_\infty=\prod_{j=1}^{N-1}(\zeta_N^je^{-2x-jt};P)_\infty\times(j=N\text{ term})$
(direct analogue of the $N=4$ case's $\prod_{j=1}^4(i^je^{-jt};P)_\infty$, with $\zeta_N=i$ read off
at $N=4$) gives resonance constant $c_j=\zeta_N^j=\zeta_6^j$ for the $j$-th factor, and that factor's
resonance locus is $\operatorname{Im}(x)=\arg(\zeta_6^j)/2\pmod\pi$ --- matching Room 8's list with
$i=j$. The leg through $x_{-1}$ sits at $\operatorname{Im}(x)=+\pi/6=\arg(\zeta_6^1)/2$, so
\[
   c_1=\zeta_6=e^{i\pi/3}
\]
is the relevant resonance constant on that leg; by conjugate symmetry the leg through $x_{-3}$
($\operatorname{Im}x=-\pi/6$) uses $c_5=\zeta_6^5=e^{-i\pi/3}=\overline{\zeta_6}$.

\section*{3. Computation along the actual leg}

Room 11's exact $N=6$ saddle data (re-verified here in 30-digit mpmath, \texttt{perl -e 'alarm 200;
exec @ARGV'} wrapper, negligible memory/disk):
\[
   x_{-1}=0.3491187578768502+0.5235987755982989\,i\ \ (\operatorname{Im}=\pi/6\text{ exactly}).
\]
Following the leg exactly as specified (constant height $\operatorname{Im}x=\pi/6$, $\operatorname{Re}x$
increasing outward from the saddle, $x(\rho)=x_{-1}+\rho$, $\rho\ge0$ real), with
$w(\rho)=c_1e^{-2x(\rho)}$:
\[
   w(\rho)=e^{i\pi/3}e^{-2x_{-1}}e^{-2\rho}
          =e^{i\pi/3}\cdot e^{-2\operatorname{Re}(x_{-1})}e^{-i\pi/3}\cdot e^{-2\rho}
          =e^{-2\operatorname{Re}(x_{-1})}\,e^{-2\rho},
\]
using $\operatorname{Im}(x_{-1})=\pi/6$ exactly, so $e^{i\pi/3}e^{-2i\operatorname{Im}(x_{-1})}
=e^{i\pi/3}e^{-i\pi/3}=1$ \textbf{exactly, algebraically, not just numerically}. So $w(\rho)$ is
\emph{real and positive} for every $\rho\ge0$: the phase cancellation is exact, confirming
$\theta^\ast=0$ (no rotation anywhere on this leg, as the papers already record), but the modulus
tells the whole story on its own.

Numerically (mpmath, 30 digits, agrees with the algebraic identity to the last printed digit,
residual imaginary part $\sim10^{-32}$ is float noise):
\[
\begin{array}{c|c|c}
\rho & w(\rho) & |1-w(\rho)|\\\hline
0.0 & 0.497461 & 0.502539\\
0.1 & 0.407287 & 0.592713\\
0.5 & 0.183006 & 0.816994\\
1.0 & 0.067324 & 0.932676\\
2.0 & 0.009111 & 0.990889
\end{array}
\]
$w(\rho)=e^{-2\operatorname{Re}(x_{-1})}e^{-2\rho}$ decreases \emph{monotonically} from
$w(0)=e^{-2\cdot0.3491188\ldots}=0.4974613\ldots$ down to $0$ as $\rho\to\infty$: the point never
approaches $w=1$ at all, it moves steadily away from it. Consequently
\[
   |1-w(\rho)|=1-w(\rho)
\]
is monotonically \emph{increasing} in $\rho$, with minimum at the saddle itself,
\[
   \boxed{|1-w(\rho)|\ \ge\ 1-e^{-2\operatorname{Re}(x_{-1})}=0.5025387\ldots>0\quad\text{for all }\rho\ge0.}
\]
By the conjugate-symmetric computation, the identical bound $|1-w(\rho)|\ge0.5025387\ldots$ holds on
the $x_{-3}$ leg with $c_5=\overline{\zeta_6}$ (verified numerically above, same values to the
printed digits, imaginary parts of $w$ again $\sim10^{-31}$ float noise).

\section*{4. Verdict}

\textbf{MODULUS-ONLY BOUND WORKS on both resonance-pinned constant-height legs, precisely stated
as follows.} Along the $N=6$ contour's constant-height legs through $x_{-1}$
($\operatorname{Im}x=\pi/6$) and $x_{-3}$ ($\operatorname{Im}x=-\pi/6$), parametrised by
$x(\rho)=x_{\mp1\text{ or }-3}+\rho$, $\rho\ge0$, and using the correctly-identified resonance
constants $c_1=\zeta_6$, $c_5=\overline{\zeta_6}$:
\[
   |1-w(\rho)|\ \ge\ 1-e^{-2\operatorname{Re}(x_{-1})}\ =\ 0.5025387\ldots\ >0,\qquad\text{for all }\rho\ge0,
\]
with equality at $\rho=0$ (the saddle) and the bound strictly improving as $\rho$ grows. The
mechanism is not a rotation argument (there is none available, $\theta^\ast=0$ exactly, confirming
the earlier diagnosis) but a pure accident of position: the saddle's real part,
$\operatorname{Re}(x_{-1})=0.3491188\ldots>0$, is large enough that $|c_1e^{-2x_{-1}}|=
e^{-2\operatorname{Re}(x_{-1})}\approx0.497<1$ starts (and stays) strictly inside the unit circle, so
the leg never gets near the singular point $w=1$ regardless of the (exactly zero) argument. This
is genuine, if narrow, progress: it replaces the ray lemma (which needs $\sin\theta\ne0$ and gives
nothing here) with an elementary, fully rigorous, explicit bound valid along the \emph{entire}
infinite leg --- not just asymptotically for $\operatorname{Re}x$ large (where part (iii)'s bound
already applied), but starting exactly at the saddle, closing precisely the gap Room 11 identified
(``it is precisely there\ldots that these two legs begin and run for a while before
$\operatorname{Re}x$ grows enough for part (iii) to take over'').

\textbf{Caveats, stated honestly.} (1) This bound is for the \emph{leading-order} $w=c_je^{-2x}$;
the lemma's actual quantities carry $O(t)$ shifts ($jt$ in the exponent, and the sum over $m$ in
$4tm$/$P$-powers) that were dropped here. Those shifts vanish as $t\to0$, which is the regime the
whole pole-family construction operates in (Room 11's calibration: $t\in\overline{D_j}$,
$|t|\le1.15\cdot10^{-3}$), and the margin found here ($0.50$, versus a bound that only needs to stay
positive) is generous enough to absorb a genuinely $O(t)$ perturbation of $\arg w$ away from exactly
$0$ --- but this was not verified quantitatively against the full $t$-dependent formula, only
argued qualitatively. A complete treatment would need to redo Part (ii)'s actual sum (the full
$\vartheta$/Poisson-summation expression, not just the single ray-lemma factor) with $t\ne0$ along
the leg, which is beyond this room's scope. (2) The bound only controls \emph{this one factor}
$(P/z_j;P)_\infty$ (equivalently the single resonance-constant piece); Part (ii)'s full bound also
needs $M'(x)$, $K'(x)$ and the $\vartheta$-sum estimate, none of which were touched here. (3) The
"$c_j=\zeta_N^j$" $N$-generic decomposition used in \S2 was reconstructed by analogy with the
$N=4$ case's stated formula, not read off an $N$-generic lemma statement verbatim (no such
statement was found in \texttt{P1f\_lemma.tex} in the material read for this room); it is internally
consistent with Room 8's independently-derived resonance loci (exact match at $i=j=1,5$), which is
good evidence it is the right generalisation, but it has not been cross-checked against a written
$N$-generic version of \texttt{lem:qi-bounds} itself, because none appears to exist yet.

\textbf{Net effect on the $N=6$ program.} This does not close $N=6$: it supplies exactly the one
missing piece Room 11 flagged as absent (a modulus-only lower bound on $|1-w|$ along the
resonance-pinned legs, near the saddle where part (iii) does not yet apply), for the
\emph{leading-order} object, with an explicit constant $0.5025\ldots$. What remains before this can
be called a closed sub-lemma: extending it to the true $t\ne0$ quantity (caveat 1), and folding it
back into the full Part (ii)-style estimate together with $M'$, $K'$, and the $\vartheta$-sum
piece (caveat 2), which is a real amount of remaining work, not a formality.

\end{document}


\section*{Correction (post-review, Reviewer AY): this does NOT constitute progress on route (b)}

The claim above is \textbf{withdrawn as stated}. Reviewer AY found a fatal, subtle gap beyond this
note's own caveats: the ray lemma's separation constant is
$d(x)=\min_{c^4=1}\mathcal D(2\operatorname{Re}x,\operatorname{dist}(\arg c-2\operatorname{Im}x,2\pi\Z))$,
$\mathcal D(a',\varphi)=\inf_{\rho\ge0}\max\{\sigma(\varphi-\rho\sin\theta),1-e^{-a'-\rho\cos\theta}\}$,
where $\rho$ is the lemma's own \emph{internal} auxiliary variable (tied to the fixed direction
$\theta=\arg t$ used in the Euler--Maclaurin sum over $m$), \emph{not} a position along the
geometric $x$-contour. This note set $\rho$ to displacement along the leg
($x(\rho)=x_{-1}+\rho$) and evaluated a single slice of the max-expression -- but $\mathcal D$ is
an \emph{infimum} over its own $\rho$, so any single evaluation (including this note's $\rho=0$
slice) is only an \emph{upper} bound on $\mathcal D$, never a lower bound, unless $\sin\theta=0$
is separately established for the construction's actual fixed $\theta$. This note's claim that
``$\theta^\ast=0$, no rotation anywhere on the leg'' conflates a different, trivially-true fact
(the phase of $w$ is constant as $x$ moves along the leg, since the leg IS the resonance locus by
construction) with the lemma's own rotation parameter, which was never checked. \textbf{The boxed
quantity $0.5025$ is a numerically correct fact about a simplified proxy, not a valid lower bound
on the ray lemma's actual separation constant.} Route (b) remains unresolved.
